\chapter{Security, Capabilities, and Microkernel Hardening}
\label{chap:security}

In traditional operating systems, security has historically been treated as an afterthought—an auxiliary layer of access control lists (ACLs) patched onto a fundamentally vulnerable monolithic foundation.

Experience has repeatedly demonstrated the perils of this model:
\begin{enumerate}
    \item \textbf{The Ambient Authority Trap:} Under standard UNIX discretionary access control (DAC), an application running as a user possesses all the privileges of that user. When a text editor or web browser is compromised, the attacker inherits the ambient authority to read SSH keys, corrupt documents, or open network sockets.
    \item \textbf{The Confused Deputy Problem:} A privileged service (such as a compiler or print daemon) can be tricked into abusing its authority on behalf of an unprivileged client because the identity of the caller is separated from the permission to access the resource.
    \item \textbf{Kernel-Space Memory Exploits:} In a monolithic C kernel, a single buffer overflow in a file system driver or network stack yields unrestricted Ring~0 code execution, bypassing all user-space security controls.
\end{enumerate}

In this chapter, we explore how Hydra achieves defense-in-depth through **Capability-Based Security**, strict microkernel isolation, and compiler-verified memory safety in Rust.

% ==============================================================================
\section{The Capability-Based Security Model}
\label{sec:capability_model}
% ==============================================================================

In a capability-based operating system, a **capability** is an unforgeable token that designates both a specific system resource (an address space, an IPC endpoint, an I/O device, or a storage dataset) and a set of permitted operations (Read, Write, Grant, Revoke):

\begin{figure}[htbp]
\centering
\begin{tikzpicture}[scale=0.9, every node/.style={transform shape}]
    \draw[fill=hydralight, draw=hydrablue, line width=1.2pt, rounded corners] (0, 0) rectangle (9, 2);
    \node[font=\bfseries\sffamily\color{hydrablue}] at (4.5, 1.5) {Hydra Capability Descriptor};
    
    \node[draw=rustkeyword, fill=white, rounded corners, minimum width=2.2cm, minimum height=0.7cm] at (1.8, 0.7) {\footnotesize Object ID};
    \node[draw=rustred, fill=white, rounded corners, minimum width=2.2cm, minimum height=0.7cm] at (4.5, 0.7) {\footnotesize Permitted Rights};
    \node[draw=hydragold, fill=white, rounded corners, minimum width=2.2cm, minimum height=0.7cm] at (7.2, 0.7) {\footnotesize Generation Tag};
\end{tikzpicture}
\caption{Structure of an Unforgeable Capability Descriptor in Hydra.}
\label{fig:cap_descriptor}
\end{figure}

Unlike ambient authority:
\begin{itemize}
    \item A process cannot access any resource unless it holds an explicit capability for it in its private **Capability Table** (\texttt{CSpace}).
    \item Capabilities cannot be forged or guessed by user processes. They are created and manipulated solely by the microkernel supervisor via type-safe handle references.
    \item When a process creates a child task or communicates with a server, it delegates only the minimal subset of capabilities required for the sub-task.
\end{itemize}

% ==============================================================================
\section{Least Privilege Across Ring~3 Servers}
\label{sec:least_privilege}
% ==============================================================================

The principle of least privilege states that every module must be able to access only the information and resources necessary for its legitimate purpose:

\begin{table}[htbp]
\centering
\small
\begin{tabular}{llll}
\toprule
\textbf{Process Name} & \textbf{I/O Port Rights} & \textbf{IPC Rights} & \textbf{Memory Rights} \\
\midrule
\texttt{ata\_driver} & Ports \texttt{0x1F0}--\texttt{0x1F7}, \texttt{0x3F6} & Server on \texttt{0x21} & DMA buffer frames only \\
\texttt{zfs\_server} & None (No I/O ports) & Server on \texttt{0x20}, Client on \texttt{0x21} & ZFS cache frames only \\
\texttt{net\_driver} & None (MMIO mapped) & Server on \texttt{0x22} & E1000 DMA rings only \\
\texttt{display\_server} & VGA ports & Server on \texttt{0x23} & Framebuffer mapping only \\
User Applications & None & Client endpoints & Private virtual address space \\
\bottomrule
\end{tabular}
\caption{Capability Allocations for Hydra System Services.}
\label{tab:cap_alloc}
\end{table}

Notice the defensive boundaries:
\begin{itemize}
    \item If an attacker compromises the OpenZFS server, they cannot send malicious packets over the network or access display memory, because the ZFS server has zero capabilities for Endpoint \texttt{0x22} (Network) or Endpoint \texttt{0x23} (Display).
    \item If a network packet triggers a flaw in the network driver, the driver cannot write to disk, because it possesses no capabilities for the ATA controller or ZFS pool.
\end{itemize}

% ==============================================================================
\section{Memory Safety as a Foundation for Kernel Security}
\label{sec:rust_security}
% ==============================================================================

More than four decades of operating system development demonstrate that human diligence cannot eliminate memory corruption bugs from large C/C++ codebases. According to security reports from Microsoft, Google Chromium, and Apple, approximately 70\% of all patched Common Vulnerabilities and Exposures (CVEs) are memory-safety violations.

Hydra eliminates entire vulnerability classes by construction:
\begin{enumerate}
    \item \textbf{Buffer Overflows and Out-of-Bounds Indexing:} Array indexing in Rust is bounds-checked at runtime, or optimized into provably safe iterator loops at compile time.
    \item \textbf{Use-After-Free and Double Free:} Rust's borrow checker enforces that memory is deallocated when its owning variable drops out of scope. Dangling pointers cannot compile.
    \item \textbf{Data Races:} Concurrent access across CPU cores requires synchronization primitives (\texttt{Mutex<T>}, \texttt{SpinLock<T>}, \texttt{RwLock<T>}) that implement the \texttt{Send} and \texttt{Sync} traits. A thread cannot access shared data without holding the lock.
\end{enumerate}

% ==============================================================================
\section{Hardware Vulnerability Mitigation}
\label{sec:hardware_mitigations}
% ==============================================================================

Modern x86\_64 processors are subject to microarchitectural speculative execution vulnerabilities (e.g., Meltdown and Spectre). Hydra incorporates modern hardening strategies:
\begin{itemize}
    \item \textbf{Kernel Page Table Isolation (KPTI):} When executing in Ring~3 user mode, the user-space PML4 page table maps only the minimal trampoline code required to handle entry into the kernel, preventing speculative reading of privileged kernel memory across privilege boundaries.
    \item \textbf{User-Space Pointer Sanitization:} All pointers passed from user space across the system call interface are validated against canonical range boundaries and checked for alignment before dereferencing.
\end{itemize}

% ==============================================================================
\section{Summary}
\label{sec:sec_summary}
% ==============================================================================

In this chapter, we analyzed Hydra's security architecture:
\begin{itemize}
    \item Capability descriptors replace ambient authority with fine-grained, unforgeable tokens of authority.
    \item The principle of least privilege strictly confines Ring~3 driver processes to minimal I/O ports and IPC endpoints.
    \item Writing the microkernel in Rust mathematically eliminates buffer overflows, use-after-free, and data races from over 95\% of the codebase.
    \item Microarchitectural mitigations protect against speculative execution vulnerabilities.
\end{itemize}

